% =====================================================================
%  SPLICE: "Where the chain goes next" -- forward-looking coda for
%  sec:wigners-friend, pointing to sec:sedenions-generations.
%  Created 2026-08-27 per JS.
%
%  THE POINT (JS, 2026-08-27). The Wigner's-friend charts stop at
%  Wigner because the Hopf tower runs out at the octonions, while the
%  text says the chain of observers obviously continues. The sedenions
%  show what the continuation adds: the observer of Wigner (fourth
%  wire, C (x) S) brings a NEW piece of structure -- the finite S_3 of
%  Brown's theorem, the generation label -- and the observer of the
%  observer of Wigner (fifth wire, the trigintaduonions) brings nothing
%  new in kind. This coda records that, so the section's termination at
%  Wigner reads as a fact about the tower rather than an economy of the
%  exposition.
%
%  WHAT THE ALGEBRA ACTUALLY SAYS (verified 2026-08-27, notebook T --
%  notebook_T_trigintaduonion_automorphisms_2026-08-27.py):
%    - Aut(S) = G_2 x S_3 (Brown 1967), DIRECT product; the new S_3 is
%      generated by the doubling involution and Brown's explicit
%      order-three map psi, which fixes 1 and the doubling unit and
%      rotates every (a, a u) plane, a imaginary, by a third of a turn.
%    - psi permutes the three octonion HALVINGS O, psi O, psi^2 O of S
%      (each doubled by the same unit; meeting only in R). It does NOT
%      permute the three copies through the quaternion core -- those are
%      cycled by a lifted G_2 element (the colour-line three-cycle); see
%      the sedenion section's two-triples remark and the corrections
%      ledger of the same date.
%    - Brown: Aut(A_{n-1}) x S_3 = Aut(A_n) with EQUALITY for n = 4, 5, 6
%      (as quoted in gresnigtgourlayvarma2023three, sec. 4.3); Eakin and
%      Sathaye prove Brown's conjecture in general (VERIFY-CITE for the
%      exact general form -- their step is Aut(A_{n+1}) = Aut(A_n) x G,
%      G in {Z/2, S_3}; over R with the standard parameters the branch is
%      S_3 at every rung >= 4, which is what notebook T checks directly
%      at rungs 4, 5, 6). Derivation algebra g_2 from the octonions on
%      (Schafer 1954; notebook T: dim Der = 14 at n = 3, 4, 5).
%  So "nothing new at the fifth wire" is stated PRECISELY: the fifth
%  doubling adds another finite S_3 and nothing else -- no connected
%  symmetry, no fibration; the increment is the same at every rung
%  thereafter. Not "silent": REPETITIVE. The prose says so.
%
%  STATUS: the algebra is ESTABLISHED (Brown; Eakin--Sathaye; Schafer)
%  and directly checked (notebook T). The reading of the repeated finite
%  factor as what the redundancy ladder consumes is a RHYME, flagged.
%  The reading of the fourth rung's discrete orientation as the
%  generation label is the sedenion section's CANDIDATE identification,
%  pointed to, not re-argued here.
%
%  PLACEMENT: a new \subsection at the END of sec:wigners-friend, after
%  the "Two cautions" paragraph. Because that paragraph reads "Two
%  cautions close the section.", the patch also changes that one word:
%  "close the section" -> "close the argument". Both anchors
%  uniqueness-asserted (the July lesson: anchors that also occur in
%  comments must be asserted unique).
%
%  CROSS-REFS OUT (all exist in the synced source except the two
%  MAP-LABELs):
%    sec:hopf-justification   -- the "two climbs" paragraph (length vs
%                                termination); this coda sharpens it
%    sec:wf-redundancy         -- the redundancy ladder (this file)
%    sec:wf-phase-promotion    -- phase promotion (this file)
%    sec:sedenions-generations -- the sedenion section (2026-08-27 rev.)
%    prop:termination          -- Hurwitz/Adams termination
%  MAP-LABELs:
%    app:notebook-record       -- the notebook-record appendix
%                                (appendices/appendix_notebook_record_
%                                2026-07-11.tex); substitute its label.
%  NEW KEYS (RIS delivered: cayley_dickson_automorphism_refs_2026-08-27.ris):
%    eakinsathaye1990automorphisms, schafer1954cayleydickson
%  IN-LIBRARY: brown1967generalized.
%
%  OPTIONAL CONSUMER-POINTER (not patched; one clause for the "two
%  climbs" paragraph of hopf_justification_section_2026-06-08.tex, after
%  "that chain has no top --- it is the bump in the carpet, unbounded in
%  length."):
%      "(unbounded in length, and --- as \S\ref{sec:wf-chain-type}
%      records --- bounded in type: after the fourth observer nothing
%      new in kind arrives)"
%
%  EDITED 2026-08-30 (per JS): one sentence added to the fourth-wire
%  paragraph anchoring the flavour sector's irremovable phase (the
%  CP-phase recognition) on this rung, with the development at
%  sec:sedenions-mixing (2026-08-30 revision). No new keys here.
%  British English; pure ASCII; \autocite; no bullets; spaced dashes.
% =====================================================================
% SPLICE-BEGIN wf-chain-type 2026-08-27

\subsection{Where the chain goes next: one more rung, then copies}
\label{sec:wf-chain-type}

The regress just described is unbounded in length, and nothing here shortens
it: the observer of Wigner carries a gauge circle of his own, and so does the
observer of that observer, as \S\ref{sec:hopf-justification} insists when it
separates the endlessness of the chain from the termination of the tower. But
the algebra says how much of the endless chain adds structure of a new kind,
and the answer is exact: one more step, and then none.

The count runs as follows. The triple system, friend, Wigner is three wires,
and three wires are the three Hopf rungs $\mathbb{C}$, $\mathbb{H}$,
$\mathbb{O}$ --- the completed link of the earlier caveat --- which is why the
charts above stop at Wigner: the chart of whoever looks in on the triple,
$S^{15}\to S^{8}$, is the last fibration the division algebras supply
(Proposition~\ref{prop:termination}). The observer of Wigner is a fourth wire,
and the fourth wire is the sedenions,
$(\mathbb{C}^{2})^{\otimes 4} = \mathbb{C}\otimes\mathbb{S}$
(\S\ref{sec:sedenions-generations}). No fibration comes with it --- the
sedenions have zero divisors and admit no map $S^{31}\to S^{16}$ --- so the
fourth observer gains no recoverable fibre of the tower's kind: the continuous
promotion of relative phase into a larger fibre, the mechanism of the previous
subsection, has no fourth instance. What the doubling adds instead is finite.
Brown's theorem gives $\mathrm{Aut}(\mathbb{S}) = G_{2}\times S_{3}$
\autocite{brown1967generalized}: the identity component is unchanged from the
octonions, and the entire gain is a component group, generated by an
order-three automorphism that fixes the new doubling direction and rotates
every imaginary direction of the completed link into its copy by a third of a
turn. The relative orientation of the link and its copy, which at the Hopf
rungs was a continuous phase promoted into a fibre, survives at the fourth rung
only as a discrete choice among three, and that choice is the generation label
(\S\ref{sec:sedenions-generations}). The fourth observer, then, sees no new
force and no new fibre; he sees the same link three ways, and cannot gauge the
difference, because a component group has no Lie algebra. His rung is also where
the flavour sector's one irremovable phase finds its anchor: three generations
are the threshold at which the mixing matrix first carries a phase no convention
removes, and \S\ref{sec:sedenions-mixing} reads that phase --- through the
paper's negativity--contextuality chain --- as this wire's contextual cargo,
real only from one rung above.

The observer of the observer of Wigner is a fifth wire, the trigintaduonions,
and here the tower has nothing new to say. Brown proved that the next two
doublings each add exactly one further $S_{3}$ and nothing else, Eakin and
Sathaye that Brown's pattern continues at every rung, and the derivation
algebra --- the Lie algebra of the identity component --- is $\mathfrak{g}_{2}$
from the octonions onward
\autocite{brown1967generalized,eakinsathaye1990automorphisms,schafer1954cayleydickson};
the fifth and sixth rungs are also checked directly in the notebook record
(Appendix~\ref{app:notebook-record}). So the fifth observer's doubling repeats
the fourth's gift --- another third-of-a-turn symmetry between the four-wire
structure and its copy, commuting with everything below --- and adds no
connected symmetry and no fibration. Nothing new in kind arrives at the fifth
wire or at any wire after it. The increments of the tower are, in order, a
force, a force, a label, and then the same finite permutation of copies at
every step. That is the precise sense in which the chain ``obviously
continues'': it continues, and after the fourth observer it has nothing further
to announce. The framework reads the repeated factor as what the redundancy
ladder of \S\ref{sec:wf-redundancy} consumes --- a permutation symmetry among
copies, which is the freedom a record has and a type does not --- and records
that reading as a structural resonance rather than a derivation.

The two climbs of \S\ref{sec:hopf-justification} are thereby sharpened. The
chain of observers is unbounded in length and bounded in type: three forces by
the third wire, one family label by the fourth, copies thereafter. The observer
of Wigner is the last observer whose arrival changes what there is.

% SPLICE-END wf-chain-type 2026-08-27
